%============================================================
% Chapter 14: Linear Basis over GF(2) (XOR Basis)
% STATUS: PLANNED  (generated by tools/outline.py; replace this file when the chapter is written)
%============================================================
\chapter{Linear Basis over GF(2) (XOR Basis)}
\label{ch:linear_basis}

\begin{tcolorbox}[enhanced, colback=white, colframe=secondaryblue, boxrule=1pt, arc=2mm]
\textbf{\large\textcolor{primaryblue}{Every set of integers spans a vector space over GF(2): a basis of at most 60 numbers answers every XOR-subset question.}}
\end{tcolorbox}

\chapterfacts{1900}{2700}{5}{10}{30}

\begin{plannedbox}[Planned Contents --- this chapter is not written yet]
\textbf{Part A --- Lessons}
\begin{enumerate}[leftmargin=*, label=\textbf{A\arabic*.}, itemsep=1pt]
    \item Integers as vectors over GF(2); span, independence and rank
    \item Inserting into a basis by Gaussian reduction from the top bit; the reduced (row echelon) form
    \item Maximum XOR of a subset, $k$-th smallest XOR, and whether $x$ is representable
    \item Counting subsets with a given XOR: $2^{n - r}$
    \item Prefix bases with positions (the latest-index trick) for range queries; merging bases
\end{enumerate}
\textbf{Part B --- Patterns} (with candidate problems, checked when the chapter is written)
\begin{enumerate}[leftmargin=*, label=\textbf{P\arabic*.}, itemsep=1pt]
    \item \textbf{Maximum and representable XOR} --- Mahmoud and Ehab and yet another xor task {\small\color{darkgray}(CF 959F)}
    \item \textbf{Rank and counting} --- (Zero XOR Subset)-less {\small\color{darkgray}(CF 1101G)}
    \item \textbf{Bases of prefixes and ranges} --- Ivan and Burgers {\small\color{darkgray}(CF 1100F)}
    \item \textbf{Merging bases}
    \item \textbf{k-th smallest XOR}
\end{enumerate}
\textbf{Connections.} XOR paths and the cycle space of a graph are the graph book; this chapter owns the basis itself.
\end{plannedbox}
